\cite{iesan1987saintvenants}.  \cite{synge1945problem}. 



\section{Less Less Old}

\subsection*{Synopsis}

We consider torsion–free flexure of a straight prismatic cylinder in 3D linearized elastostatics. The axis is the unit vector $\mathbf{i}$; the cross–section is $\Omega\subset \mathbf{i}^\perp$ with outward unit normal $\bm{\nu}$; points are $x\,\mathbf{i}+\bm{r}$ with $\bm{r}=(0,y,z)$. 
The material is homogeneous, isotropic, with Lamé moduli $(\lambda,G)$.

\emph{Saint–Venant class.} On a prismatic cylinder with traction–free lateral surface and loads prescribed only by end \emph{resultants}, the stress admits the representation
\begin{equation}
\bm{T}=\sigma\,\mathbf{i}\otimes\mathbf{i}\;+\;\bm{\tau}\otimes\mathbf{i}\;+\;\mathbf{i}\otimes\bm{\tau},
\qquad
\sigma:=\mathbf{i}\cdot\bm{T}\mathbf{i},\quad
\bm{\tau}:=(\mathbf{i}_\beta\otimes \mathbf{i}_\beta)\bm{T}\mathbf{i}\in \mathbf{i}^\perp,
\label{eq:frame-stress-tensor}
\end{equation}
and the side condition $\bm{T}\bm{n}=\mathbf{0}$ (with $\bm{n}=\bm{\nu}\in\mathbf{i}^\perp$) reduces to $\bm{\tau}\cdot\bm{\nu}=0$ on $\partial\Omega$. Symmetry of $\bm{T}$ and the vanishing in–plane block justify \eqref{eq:frame-stress-tensor}.

\emph{Canonical selection.} The relaxed formulation (resultants only) is not unique \emph{a priori}. 
We select the canonical Saint–Venant flexure (Ieşan): stresses depend at most linearly on $x$ (axial coordinate), torsional amplitude vanishes, and the cross–sectional problem reduces to a single scalar warping function with Neumann data.

\subsection*{Problem}

\paragraph{Body and fields.}
Prismatic cylinder $\mathcal{C}=\{\,x\,\mathbf{i}+\bm{r}\ :\ x\in(0,L),\ \bm{r}\in\Omega\subset\mathbf{i}^\perp\,\}$, with unknown single–valued displacement $\bm{u}$ and stress $\bm{T}$ satisfying
\[
\operatorname{div}\bm{T}=\mathbf{0},\qquad
\bm{T}=\lambda(\operatorname{tr}\bm{\varepsilon})\mathbf{I}+2G\,\bm{\varepsilon},\qquad
\bm{\varepsilon}=\tfrac12(\nabla \bm{u}+\nabla \bm{u}^{\!\top}).
\]

\paragraph{Boundary and loading.}
Lateral surface $(0,L)\times\partial\Omega$ is traction–free: $\bm{T}\bm{n}=\mathbf{0}$. At $x=L$ prescribe a transverse resultant force $\bm{F}=F\,\bm{f}$ with $\bm{f}\in\mathbf{i}^\perp$ and zero resultant couple (torsion–free). The end $x=0$ is clamped. No body forces.

\paragraph{Saint–Venant stress hypothesis.}
Restrict to the class \eqref{eq:frame-stress-tensor} with $\bm{\tau}\cdot\bm{\nu}=0$ on $\partial\Omega$.

\paragraph{Canonical flexure postulate (Ieşan).}
Among relaxed solutions, select the Saint–Venant flexure whose stresses depend at most linearly on $x$; equivalently, the torsional amplitude is zero and the cross–sectional remainder is governed by a scalar warping function.

These postulates determine a unique solution up to rigid motions.

\subsection*{Solution}

\subsubsection*{Resultants and axial reduction}

With \eqref{eq:frame-stress-tensor}, equilibrium splits as
\[
\partial_x \bm{\tau}=\mathbf{0},\qquad
\partial_x \sigma+\operatorname{div}_{\bm{r}}\bm{\tau}=0.
\]
Hence $\bm{\tau}=\bm{\tau}(\bm{r})$ and
\begin{equation}
\sigma(x,\bm{r})=(L-x)\,q(\bm{r}),
\qquad
q(\bm{r}):=\operatorname{div}_{\bm{r}}\bm{\tau}(\bm{r}).
\label{eq:sigma-q}
\end{equation}
The shear and bending resultants are
\[
\bm{V}=\int_\Omega \bm{\tau}\,dA=\bm{F},\qquad
\bm{M}(x)=\int_\Omega \bm{r}\times\big(\sigma\,\mathbf{i}\big)\,dA=(L-x)\,\mathbf{i}\times\bm{F}.
\]

\subsubsection*{Section inertias and load amplitudes}

Let the (planar) second–moment operator $\bm{\mathsf{I}}:\mathbf{i}^\perp\to \mathbf{i}^\perp$ be defined in the basis $\{\mathbf{i}_y,\mathbf{i}_z\}$ by the symmetric matrix
\[
[\bm{\mathsf{I}}]=
\begin{pmatrix}
I_y & -I_{yz}\\
-I_{yz} & I_z
\end{pmatrix},
\qquad
I_y=\int_\Omega z^2\,dA,\ \ I_z=\int_\Omega y^2\,dA,\ \ I_{yz}=\int_\Omega y z\,dA,
\]
with the origin at the area centroid (so $\int_\Omega \bm{r}\,dA=\mathbf{0}$). The canonical flexure amplitude $\bm{b}\in\mathbf{i}^\perp$ is fixed by the resultant relation
\begin{equation}
\bm{F}=E\,\bm{\mathsf{I}}\,\bm{b},
\label{eq:F-EIb}
\end{equation}
and a scalar $b_3\in\mathbb{R}$ (axial-shift amplitude) determined by $\int_\Omega q\,dA=0$; with centroidal origin one has $b_3=0$.

\subsubsection*{Warping problem and stresses}

Introduce a scalar warping function $\varphi:\Omega\to\mathbb{R}$ and the quadratic, divergence–free polynomial field
\begin{equation}
\bm{\Pi}(\bm{r};\bm{b})
=
\Big[-(\mathbf{i}_y\cdot\bm{b})\,y z-\tfrac12(\mathbf{i}_z\cdot\bm{b})\,(y^2-z^2)\Big]\mathbf{i}_y
+\Big[\tfrac12(\mathbf{i}_y\cdot\bm{b})\,(y^2-z^2)+(\mathbf{i}_z\cdot\bm{b})\,y z\Big]\mathbf{i}_z,
\qquad \operatorname{div}_{\!\bm{r}}\bm{\Pi}\equiv 0.
\label{eq:Pidef}
\end{equation}
The canonical Saint–Venant flexure is obtained by solving the \emph{scalar} Poisson–Neumann problem
\begin{equation}
\Delta_{\!\bm{r}}\varphi = -\,2\big(\bm{b}\cdot\bm{r}+b_3\big)\quad\text{in }\Omega,
\qquad
\partial_{\bm{\nu}}\varphi = -\,\nu\,\bm{\Pi}(\bm{r};\bm{b})\cdot\bm{\nu}\quad\text{on }\partial\Omega,
\qquad
\int_\Omega \varphi\,dA=0,
\label{eq:phiBVP}
\end{equation}
and setting the shear and axial stresses to
\begin{equation}
\boxed{\ \bm{\tau}(\bm{r})=2G\Big[\nabla_{\!\bm{r}}\varphi(\bm{r})+\nu\,\bm{\Pi}(\bm{r};\bm{b})\Big]\ },
\qquad
\boxed{\ \sigma(x,\bm{r})=(L-x)\,\operatorname{div}_{\!\bm{r}}\bm{\tau}(\bm{r})\ }.
\label{eq:stress-final}
\end{equation}
By construction, $\bm{\tau}\cdot\bm{\nu}=0$ on $\partial\Omega$ and \eqref{eq:sigma-q} holds with
\[
\operatorname{div}_{\!\bm{r}}\bm{\tau}=2G\,\Delta_{\!\bm{r}}\varphi=-\,4G\big(\bm{b}\cdot\bm{r}+b_3\big),
\]
so that $\int_\Omega \sigma(x,\bm{r})\,dA=0$ for all $x$. The resultants $\bm{V}=\bm{F}$ and $\bm{M}(x)=(L-x)\,\FrameBasis\times\bm{F}$ are satisfied.

\paragraph{Displacement representation.}
A representative displacement generating \eqref{eq:stress-final} under Hooke’s law is
\begin{equation}
\bm{u}(x,\bm{r})=\bm{a}_0+\bm{a}_1\,x+(\bm{\theta}_0+\bm{\theta}_1 x)\times \bm{r}\;+\;\alpha\,\mathbf{i}\,\varphi(\bm{r}),
\label{eq:svf-disp}
\end{equation}
where $\bm{a}_0,\bm{a}_1,\bm{\theta}_0,\bm{\theta}_1$ are determined (up to an inessential rigid motion) by the clamp and end–resultant conditions, and $\alpha$ is a warping amplitude (absorbed into $\bm{b}$ by \eqref{eq:F-EIb} if desired). In the canonical family the $x$–dependence is at most linear.

\subsection*{Observations}

\paragraph{Helmholtz content on the section.}
With \eqref{eq:stress-final},
\[
\bm{\tau} = 2G\,\gradS \varphi + 2G\,\nu\,\bm{\Pi}(\bm{r};\bm{b}),
\]
is the orthogonal sum of an irrotational part and a divergence–free quadratic ``Poisson correction." Hence $\operatorname{curl}_{\!\bm{r}}\bm{\tau}=\mathbf{0}$ iff $\nu=0$; for general $\nu$ the shear is not curl–free. There is no harmonic component (origin at the centroid; displacement formulation).

\paragraph{Alternative flexure fields.}
Any self–equilibrated sectional field $\widehat{\bm{\tau}}$ with
$\operatorname{div}_{\!\bm{r}}\widehat{\bm{\tau}}=0$, $\widehat{\bm{\tau}}\cdot\bm{\nu}=0$, and vanishing force and moment resultants may be superposed without altering the end resultants; the canonical Saint–Venant selection discards such additions (minimum data, linear in $x$).

\paragraph{Summary.}
Given $\bm{F}$ and $\Omega$ (centroidal axes):
(i) solve \eqref{eq:F-EIb} for $\bm{b}$ (take $b_3=0$);
(ii) solve \eqref{eq:phiBVP} for $\varphi$ (unique up to mean zero);
(iii) form $\bm{\tau},\sigma$ from \eqref{eq:stress-final};
(iv) recover a displacement using \eqref{eq:svf-disp}.

\section{Less Old}
\subsection*{Synopsis}

We consider the flexure (transverse tip load, torsion–free) of a straight prismatic cylinder in 3D linearized elastostatics. 
Let the axis be the unit vector $\mathbf{i}$ and the cross–section be $\Omega\subset\mathbf{i}^\perp$ with outward unit normal $\bm{\nu}$; points are $x\,\mathbf{i}+\bm{r}$ with $\bm{r}=(0,y,z)$. The Cauchy stress is everywhere symmetric and the material is isotropic with Lamé moduli $(\lambda,G)$.

We work in the relaxed Saint–Venant setting: end tractions are prescribed only through their resultants, the lateral boundary is traction–free, and we seek the canonical Saint–Venant solution within this class. 
Two key facts due to Ieşan guide the formulation \citep{iesan1986saintvenants}: (i) the relaxed problem is not unique \emph{a priori}, but the canonical Saint–Venant solutions (extension, bending, torsion, flexure) can be picked out by an axial–derivative criterion; in particular, for flexure the canonical field is obtained by Ieşan’s representation (his Eq.\ 2.16) and is unique up to rigid motions, and 
(ii) the ``minimum–energy selects Saint–Venant" characterization that holds for torsion fails for flexure except at $\nu=0$—hence the need for the canonical construction. 

\smallskip
\noindent\textit{Salient observations.}
\begin{itemize}
\item Canonical flexure is characterized by Ieşan’s integral–plus–sectional representation; in particular, the torsional amplitude vanishes ($b_4=0$), and a scalar cross–sectional function solves a plane Neumann problem. 
\item Stresses in the canonical flexure depend at most \emph{linearly} on the axial coordinate. (Via the ``primary solution of Saint–Venant type" whose axial derivative is rigid.)
\item The energy–minimizer characterization that is decisive for torsion applies to flexure only if $\nu=0$; otherwise one uses the canonical representation to remove ambiguity.
\end{itemize}

\subsection*{Problem}

Let $\mathcal{C}=\{\,x\,\mathbf{i}+\bm{r}:\ x\in(0,L),\ \bm{r}\in\Omega\subset\mathbf{i}^\perp\,\}$, with $\Omega$ simply connected and of class $C^{1}$. The lateral boundary $(0,L)\times\partial\Omega$ is traction–free. At $x=L$ a resultant force $\bm{F}=F\,\bm{f}$ with $\bm{f}\in\mathbf{i}^\perp$ acts, with \emph{zero resultant couple} (torsion–free flexure). The end $x=0$ is clamped. Body forces vanish.

\paragraph{Field equations.}
\[
\operatorname{div}\bm{T}=\mathbf{0},\qquad
\bm{T}=\lambda(\operatorname{tr}\bm{\varepsilon})\mathbf{I}+2G\bm{\varepsilon},\qquad
\bm{\varepsilon}=\tfrac12(\nabla\bm{u}+\nabla\bm{u}^{\!\top}).
\]

\paragraph{Saint–Venant class and stress form.}
In the Saint–Venant class for a prismatic cylinder with free sides, all nonzero stress components have one index along $\mathbf{i}$. With the orthonormal frame $\{\mathbf{i},\mathbf{i}_y,\mathbf{i}_z\}$ spanning $\mathbb{R}^3$ and $\mathbf{i}_y,\mathbf{i}_z\in\mathbf{i}^\perp$, write
\begin{equation}
\bm{T}=\sigma\,\mathbf{i}\otimes\mathbf{i}+\bm{\tau}\otimes\mathbf{i}+\mathbf{i}\otimes\bm{\tau},
\qquad
\sigma:=\mathbf{i}\cdot\bm{T}\mathbf{i},\quad
\bm{\tau}:=(\mathbf{i}_\beta\otimes\mathbf{i}_\beta)\bm{T}\mathbf{i}\in\mathbf{i}^\perp.
\label{eq:SVform}
\end{equation}
Symmetry of $\bm{T}$ and the vanishing in–plane block, together with $\bm{T}\bm{n}=\mathbf{0}$ on the side (where $\bm{n}=\bm{\nu}\in\mathbf{i}^\perp$), justify \eqref{eq:SVform} and reduce the side traction to the scalar condition $\bm{\tau}\cdot\bm{\nu}=0$ on $\partial\Omega$.

\paragraph{Relaxed Boundary}
At $x=L$ the end traction $\bm{t}=\bm{T}\mathbf{i}$ satisfies
\[
\int_\Omega \bm{t}\,dA=\bm{F},\qquad
\int_\Omega \bm{r}\times\bm{t}\,dA=\mathbf{0}\quad\text{(torsion–free)}.
\]
We impose the \emph{canonical} Saint–Venant selection for flexure: the solution belongs to Ieşan’s class generated by his representation (Eq.\ 2.16), which fixes the amplitudes uniquely (up to rigid motions) and enforces the ``linear–in–axis" character \citep{iesan1987saintvenants}.

\subsection*{Solution}

\subsubsection*{Resultants and axial reduction}
With \eqref{eq:SVform}, equilibrium splits into
\[
\partial_x\bm{\tau}=\mathbf{0},\qquad
\partial_x\sigma+\operatorname{div}_{\!\bm{r}}\bm{\tau}=0,
\]
so $\bm{\tau}=\bm{\tau}(\bm{r})$ is sectional and
\begin{equation}
\sigma(x,\bm{r})=(L-x)\,q(\bm{r}),\qquad
q(\bm{r}):=\operatorname{div}_{\!\bm{r}}\bm{\tau}(\bm{r}).
\label{eq:sigmaq}
\end{equation}
The bending moment is $\bm{M}(x)=\displaystyle\int_\Omega \bm{r}\times(\sigma\,\mathbf{i})\,dA=(L-x)\,\mathbf{i}\times\bm{F}$, and the shear resultant is $\bm{V}=\displaystyle\int_\Omega\bm{\tau}\,dA=\bm{F}$.

\subsubsection*{Canonical cross–sectional fields (Ieşan)}
Ieşan’s construction shows that the plane elastic remainder vanishes (modulo a 2D rigid motion) and that the sectional scalar $\varphi:\Omega\to\mathbb{R}$ must solve a Poisson–Neumann problem driven by a constant vector $\bm{b}\in\mathbf{i}^\perp$ and a constant $b_3\in\mathbb{R}$:
\begin{equation}
\Delta_{\!\bm{r}}\varphi=-2(\bm{b}\cdot\bm{r}+b_3)\quad\text{in }\Omega,
\qquad
\partial_{\bm{\nu}}\varphi=\text{(linear form in $(\bm{b},\bm{r},\bm{\nu},\nu)$) on }\partial\Omega,
\label{eq:phiBVP}
\end{equation}
with the boundary datum the standard ``Poisson correction" from the free–side condition (explicit component form given in his (2.21)). The torsional amplitude vanishes, $b_4=0$. 

In principal centroidal axes, the Saint–Venant relation fixes $\bm{b}$ from the force resultant,
\begin{equation}
\bm{F}=E\,\bm{\mathsf{I}}\,\bm{b},\qquad
\bm{\mathsf{I}}:\mathbf{i}^\perp\!\to\!\mathbf{i}^\perp\ \text{the second–moment operator about }\mathbf{i},
\label{eq:SV-FI}
\end{equation}
i.e.\ the diagonal form $F_y=E I_y b_y,\ F_z=E I_z b_z$ (and the coupling $I_{yz}$ if axes are not principal).

\subsubsection*{Shear and axial stress}
The canonical Saint–Venant flexure shear decomposes as
\begin{equation}
\boxed{\ \bm{\tau}(\bm{r})=2G\,\nabla_{\!\bm{r}}\varphi(\bm{r})\;+\;2G\,\nu\,\bm{\tau}^{\!(p)}(\bm{r};\bm{b})\ },
\label{eq:tau-decomp}
\end{equation}
where $\bm{\tau}^{\!(p)}(\cdot;\bm{b})$ is a \emph{fixed polynomial} vector field on $\Omega$ (quadratic in $\bm{r}$, linear in $\bm{b}$) determined by the Poisson effect; it has zero divergence in $\Omega$ and supplies the nonzero Neumann datum in \eqref{eq:phiBVP} so that $\bm{\tau}\cdot\bm{\nu}=0$ on $\partial\Omega$. In particular,
\[
\operatorname{div}_{\!\bm{r}}\bm{\tau}=2G\,\Delta_{\!\bm{r}}\varphi=-4G\,(\bm{b}\cdot\bm{r}+b_3),
\]
and therefore, from \eqref{eq:sigmaq},
\begin{equation}
\boxed{\ \sigma(x,\bm{r})=-4G\,(L-x)\,(\bm{b}\cdot\bm{r}+b_3)\ }.
\label{eq:sigma-final}
\end{equation}
Equations \eqref{eq:phiBVP}–\eqref{eq:sigma-final} are exactly Ieşan’s canonical flexure fields, phrased in vector form; compare his (2.18)–(2.22).

\subsubsection*{Displacement with axial warping amplitude}
Let $\alpha$ be the (scalar) \emph{warping amplitude} (constant along $x$ in the canonical family; its value is absorbed by the choice of $\bm{b}$ via \eqref{eq:SV-FI}). A representative displacement that generates \eqref{eq:tau-decomp}–\eqref{eq:sigma-final} under Hooke’s law is
\begin{equation}
\bm{u}(x,\bm{r})=\bm{a}(x)+\bm{\theta}(x)\times\bm{r}\;+\;\alpha\,\mathbf{i}\,\varphi(\bm{r}),
\label{eq:svf-disp}
\end{equation}
where $\bm{a},\bm{\theta}$ are the rigid motions of the cross–section (determined up to inessential constants by the end data) and $\varphi$ solves \eqref{eq:phiBVP}. In Ieşan’s construction the ``plane elastic remainder" is null (modulo a 2D rigid motion), and the torsional amplitude is zero; the remaining $x$–dependence is at most linear and is carried by $\bm{a},\bm{\theta}$.

\subsection*{Observations}

\paragraph{Helmholtz content of $\bm{\tau}$.}
On the section $\Omega$ (simply connected), the canonical Saint–Venant shear \eqref{eq:tau-decomp} has the Helmholtz–Hodge decomposition
\[
\bm{\tau}= \underbrace{2G\,\nabla_{\!\bm{r}}\varphi}_{\text{irrotational}} \;+\; \underbrace{2G\,\nu\,\bm{\tau}^{\!(p)}(\cdot;\bm{b})}_{\text{solenoidal (divergence-free)}},
\]
with no harmonic component. Thus, \emph{in general} ($\nu\neq 0$) the shear is \emph{not} curl–free; the solenoidal piece is the familiar quadratic ``Poisson correction" that appears in classical component formulas (Love; Pilkey–Romano). When $\nu=0$, $\bm{\tau}$ collapses to a pure gradient and the energy–minimizer characterization holds.

\paragraph{Vector form vs. classical components.}
Writing \eqref{eq:tau-decomp} in the principal basis $\{\mathbf{i}_y,\mathbf{i}_z\}$ reproduces the standard expressions
\[
\bm{\tau}\cdot\mathbf{i}_y=\frac{2G}{1+\nu}\Big(\partial_y\Phi+\nu\,\Pi_y(y,z)\Big),\quad
\bm{\tau}\cdot\mathbf{i}_z=\frac{2G}{1+\nu}\Big(\partial_z\Phi+\nu\,\Pi_z(y,z)\Big),
\]
where $\Phi$ is a scaled version of $\varphi$ and $(\Pi_y,\Pi_z)$ are the quadratic polynomials in $(y,z)$ fixed by $\bm{b}$ and the section inertias (the textbook forms you quoted). The vector presentation \eqref{eq:tau-decomp} simply collects these into ``gradient $+$ solenoidal polynomial" with $\Delta\varphi$ and the Neumann datum given by \eqref{eq:phiBVP} (Ieşan’s (2.21)).

\paragraph{What makes the solution canonical.}
Ieşan’s representation (his Eq.\ 2.16) together with the axial–derivative characterization of ``primary solutions of Saint–Venant type" picks out \emph{exactly} the classical Saint–Venant flexure field (no residual ambiguity, up to rigid motions). 
This is the precise replacement for an energy–minimizer principle in flexure when $\nu\neq 0$. 

\bigskip
\noindent\textbf{Minimal recipe.}
\begin{enumerate}
\item From $\bm{F}$ and the section inertias (in principal axes), solve \eqref{eq:SV-FI} for $\bm{b}$ (and set $b_4=0$).
\item Solve the plane problem \eqref{eq:phiBVP} for $\varphi$ (Neumann data as in Ieşan (2.21)); fix $\varphi$ up to an additive constant by $\int_\Omega\varphi\,dA=0$.
\item Form $\bm{\tau}$ and $\sigma$ from \eqref{eq:tau-decomp}–\eqref{eq:sigma-final}; the side condition $\bm{\tau}\cdot\bm{\nu}=0$ holds by construction; resultants satisfy $\int_\Omega\bm{\tau}=\bm{F}$ and $\int_\Omega \bm{r}\times\bm{\tau}=\mathbf{0}$.
\item (Optional) Recover a displacement using \eqref{eq:svf-disp} (choose $\bm{a},\bm{\theta}$ to meet the clamp and end conditions; $\alpha$ can be absorbed in $\bm{b}$).
\end{enumerate}

SaintVenant's the a \emph{relaxed} statement in which the pointwise assignment of the terminal tractions 1s replaced by prescribing the corresponding resultant for ce and resultant moment.
The relaxed statement of the problem fails to characterize the solution uniquely.

Saint Venant's flexure solution is obtained if the stress field depends on tbe axial coordinate at most linearly.
Great catch — I slipped in that ``simply connected" caveat because it matters for the Hodge/Helmholtz split, and (importantly) because **Ieşan’s paper itself assumes a simply connected section** for the canonical construction. That’s explicit in his setup for the cylinder: the cross-section is taken ``a simply connected regular region." 


On a **multiply connected** section $\Omega\subset\mathbf{i}^\perp$ with boundary consisting of several components, the 2-D Hodge/Helmholtz decomposition (with the Saint-Venant side condition $\bm\tau\cdot\bm\nu=0$ on each boundary component) is
\[
\bm\tau = \nabla\phi+\nabla^{!\perp}\psi+\bm h,
\]
where

* $\nabla\phi$ is the irrotational (curl-free) part;
* $\nabla^{!\perp}\psi$ is the solenoidal (divergence-free) part;
* $\bm{h}$ is a harmonic Neumann field: $\operatorname{div}\bm h=0$, $\operatorname{curl}\bm h=0$, $\bm h\cdot\bm\nu=0$ on $\partial\Omega$.

For simply connected $\Omega$, $\bm h\equiv \mathbf{0}$. 
For multiply connected $\Omega$ with $g$ holes, $\dim{\bm h}=g$: you can add independent ``circulation loops" tangent to the boundaries without changing $\operatorname{div}\bm\tau$ or the side traction. 
This is exactly the finite-dimensional ambiguity that disappears when $\Omega$ is simply connected.

A harmonic Neumann field $\bm h$ has $\operatorname{div}\bm h=0$ and $\bm h\cdot\bm\nu=0$; by adjusting $\nabla\phi$ one can still meet the force/moment resultants. 
So $\bm h$ is *not* eliminated by the standard relaxed constraints; it remains undetermined unless you add more information.


For a simply connected section there is no harmonic subspace, so the canonical solution is unique (up to rigid motions) and has the structure we wrote (gradient warping (+) Poisson-correction polynomial).

* If you truly want **arbitrary** sections (holes allowed) and still want a *canonical* SV flexure field, you need one extra, natural condition to kill the $g$-dimensional harmonic freedom. 
Two equivalent ways that practitioners use:
\begin{enumerate}
\item **Zero circulation constraints**: for each hole boundary component $\Gamma_k$, impose
     \[
     \oint_{\Gamma_k} \bm\tau\cdot\bm t,ds = 0,\qquad k=1,\dots,g,
     \]
     where $\bm t$ is the positively oriented tangent. This sets $\bm h\equiv 0$.
\item **Minimum complementary energy within the sectional class**: minimize $\int_\Omega \tfrac{1}{4G}\lvert \bm\tau\rvert^2 \,dA$ subject to the Saint-Venant constraints (resultants, equilibrium with $\sigma$, $\bm\tau\cdot\bm\nu=0$). Because the harmonic space is orthogonal in the $L^2$ inner product, the minimizer drops $\bm{h}$.
     (This is distinct from the global 3-D energy characterization that fails for flexure at $\nu\neq 0$; here we’re only selecting among **sectional** fields $\bm\tau$ consistent with the SV class.)
\end{enumerate}

Either way, you recover exactly the same ``gradient (+) fixed polynomial Poisson-correction" structure we wrote — with **no** extra harmonic term — even when $\Omega$ has holes.

For multiply connected $\Omega$, we additionally require zero tangential circulation of $\bm\tau$ on each boundary component (equivalently, we select the minimum-complementary-energy representative).
With that single add-on, the **Helmholtz content** of the canonical SV flexure shear stays:
\[
\bm\tau = 2G\,\nabla_{\!\bm{r}}\varphi + 2G\,\nu\,\bm\tau^{(p)}(\bm{r};\bm b),
\]


\section{Old}



\subsection*{Problem}

Let $\mathcal{C}=\{\,x\,\mathbf{i}+\bm{r}\;:\;x\in(0,L),\ \bm{r}\in\Omega\subset\mathbf{i}^\perp\,\}$ be a straight prismatic cylinder with axis unit vector $\mathbf{i}$, cross‐section $\Omega$ (bounded, Lipschitz), and outward unit normal $\bm{\nu}$ on $\partial\Omega$. A constant transverse force $\bm{F}=F\,\bm{f}$ ($\bm{f}\in\mathbf{i}^\perp$) is applied on the free end $\{L\}\times\Omega$ with no end couple; the clamped end $\{0\}\times\Omega$ is kinematically fixed; the lateral surface $(0,L)\times\partial\Omega$ is traction–free. Body forces vanish. The material is linear, isotropic, with Lamé moduli $(\lambda,G)$.

We seek a static solution $(\bm{u},\bm{T})$ of 3D linearized elastostatics on $\mathcal{C}$ satisfying
\[
\operatorname{div}\bm{T}=\mathbf{0},
\qquad 
\bm{T}=\lambda(\operatorname{tr}\bm{\varepsilon})\mathbf{1}+2G\bm{\varepsilon},
\qquad 
\bm{\varepsilon}=\tfrac12(\nabla\bm{u}+\nabla\bm{u}^{\!\top}),
\]
with boundary conditions
\[
\bm{u}=\mathbf{0}\ \text{on }\{0\}\times\Omega,\qquad
\bm{T}\bm{n}=\mathbf{0}\ \text{on }(0,L)\times\partial\Omega,
\]
and end traction $\bm{t}=\bm{T}\mathbf{i}$ on $\{L\}\times\Omega$ satisfying the Saint‐Venant resultants
\[
\int_\Omega \bm{t}\,dA=\bm{F},\qquad 
\int_\Omega \bm{r}\times\bm{t}\,dA=\mathbf{0}\quad\text{(torsion free)}.
\]

\subsection*{Saint–Venant class and stress representation}

Introduce the orthonormal frame $\{\mathbf{i},\mathbf{i}_y,\mathbf{i}_z\}$ with $\mathbf{i}_y,\mathbf{i}_z\in\mathbf{i}^\perp$ and write $\bm{r}=(0,y,z)$. 
The Saint–Venant elastostatic class for prismatic cylinders assumes that, away from the ends, in‐plane normal and in‐plane shear stresses vanish, so every nonzero stress component has one index along $\mathbf{i}$. 
Equivalently,
\begin{equation}
\bm{T}
=\sigma\,\mathbf{i}\otimes\mathbf{i}\;+\;\bm{\tau}\otimes\mathbf{i}\;+\;\mathbf{i}\otimes\bm{\tau},
\qquad
\sigma:=\mathbf{i}\cdot\bm{T}\mathbf{i},\quad
\bm{\tau}:=(\mathbf{i}_\beta\otimes\mathbf{i}_\beta)\bm{T}\mathbf{i}\in\mathbf{i}^\perp,
\label{eq:frame-stress-tensor}
\end{equation}
where $\{\mathbf{i}_\beta\}$ spans $\mathbf{i}^\perp$. 
This form is \emph{justified} as follows:
(i) symmetry of $\bm{T}$ permits only three independent blocks; 
(ii) traction on the lateral surface has normal $\bm{n}=\bm{\nu}\in\mathbf{i}^\perp$, so $\bm{T}\bm{n}=(\bm{\tau}\cdot\bm{n})\,\mathbf{i}$; imposing $\bm{T}\bm{n}=\mathbf{0}$ on $(0,L)\times\partial\Omega$ forces $\bm{T}$ to have no purely in‐plane block (else $\bm{T}\bm{n}$ would have in‐plane parts), and reduces the boundary condition to
\begin{equation}
\bm{\tau}\cdot\bm{\nu}=0\quad\text{on }\partial\Omega.
\label{eq:side-BC}
\end{equation}

% \subsection*{Equations reduced on the cross–section}

With the ansatz \eqref{eq:frame-stress-tensor} and $\bm{\tau}=\bm{\tau}(\bm{r})$ (independent of $x$), equilibrium $\operatorname{div}\bm{T}=\mathbf{0}$ splits orthogonally into
\begin{equation}
\partial_x\bm{\tau}=\mathbf{0},\qquad 
\partial_x\sigma+\operatorname{div}_{\bm{r}}\bm{\tau}=0.
\label{eq:equil-split}
\end{equation}
Hence $\bm{\tau}$ is a section field, and axial stress \(\sigma\) takes the form $\sigma(x,\bm{r})=\sigma(L,\bm{r})-(x-L)\,q(\bm{r})$ with
\begin{equation}
q(\bm{r}):=\operatorname{div}_{\bm{r}}\bm{\tau}(\bm{r}).
\label{eq:q-def}
\end{equation}
Choosing the tip to carry no axial stress ($\sigma(L,\bm{r})=0$) enforces the Saint–Venant ``force without couple" end condition; all axial stress is then induced by the sectional divergence of $\bm{\tau}$ as $x$ decreases from $L$.

\subsection*{Resultants and torsion–free constraint}

Define the sectional shear resultant and bending moment (about the axis through the origin of $\Omega$) by
\[
\bm{V}:=\int_\Omega \bm{\tau}\,dA\in\mathbf{i}^\perp,\qquad
\bm{M}(x):=\int_\Omega \bm{r}\times\big(\sigma(x,\bm{r})\,\mathbf{i}\big)\,dA\in\mathbf{i}^\perp.
\]
The end conditions impose
\begin{equation}
\bm{V}=\bm{F},\qquad 
\int_\Omega \bm{r}\times\bm{\tau}\,dA=\mathbf{0}\quad(\text{torsion free}).
\label{eq:resultants}
\end{equation}
Differentiating $\bm{M}$ in $x$ and using \eqref{eq:equil-split} with a cross–sectional integration by parts yields the Saint–Venant balance of resultants
\begin{equation}
\bm{M}'(x)=\mathbf{i}\times\bm{V}=\mathbf{i}\times\bm{F},\qquad \bm{M}(L)=\mathbf{0}\ \Rightarrow\ \bm{M}(x)=(L-x)\,\mathbf{i}\times\bm{F}.
\label{eq:moment-balance}
\end{equation}

\subsection*{Shear field as a plane problem}

The shear field $\bm{\tau}:\Omega\to\mathbf{i}^\perp$ is determined solely on $\Omega$ by
\begin{equation}
\begin{aligned}
\bm{\tau}\in H(\mathrm{div};\Omega;\mathbf{i}^\perp),\\ 
\bm{\tau}\cdot\bm{\nu}&=0\ \text{on }\partial\Omega,\\
\int_\Omega \bm{\tau}\,dA&=\bm{F},\\
\int_\Omega \bm{r}\times\bm{\tau}\,dA&=\mathbf{0}.
\end{aligned}
\label{eq:tau-constraints}
\end{equation}
Among all such fields, the Saint–Venant solution is the (unique) minimizer of complementary energy
\begin{equation}
\bm{\tau}=\underset{\tilde{\bm{\tau}}\ \text{satisfying }\eqref{eq:tau-constraints}}{\arg\min}\;
\frac{1}{4G}\int_\Omega |\tilde{\bm{\tau}}|^2\,dA.
\label{eq:SV-variational}
\end{equation}
The Euler–Lagrange equations for \eqref{eq:SV-variational} yield a \emph{divergence–form} plane boundary value problem:
\begin{equation}
-\operatorname{div}_{\bm{r}}\big(\bm{\tau}\big)=q\ \text{in }\Omega,\qquad 
\bm{\tau}\cdot\bm{\nu}=0\ \text{on }\partial\Omega,
\label{eq:tau-strong}
\end{equation}
with the scalar source $q$ fixed implicitly by the two global constraints in \eqref{eq:tau-constraints} (force and zero torsional moment). 

\subsection*{Axial stress reconstruction}

With $\bm{\tau}$ known, the axial stress follows from \eqref{eq:equil-split}–\eqref{eq:q-def} as
\begin{equation}
\sigma(x,\bm{r})=(L-x)\,q(\bm{r}),\qquad 
q(\bm{r})=\operatorname{div}_{\bm{r}}\bm{\tau}(\bm{r}),
\label{eq:sigma-from-tau}
\end{equation}
and the moment balance \eqref{eq:moment-balance} is satisfied identically. By construction $\int_\Omega \sigma(x,\bm{r})\,dA=0$ for all $x$ (no axial force), because $\int_\Omega q\,dA=\int_{\partial\Omega}\bm{\tau}\cdot\bm{\nu}\,ds=0$.

\subsection*{Displacement and warping}

Let $\bm{u}$ be any displacement field whose strain produces \eqref{eq:frame-stress-tensor} under Hooke's law. 
Decompose $\bm{u}$ into a rigid cross–sectional motion and an axial warping:
\[
\bm{u}(x,\bm{r})=\bm{a}(x)+\bm{\theta}(x)\times\bm{r}\;+\;\big(\mathbf{i}\,\varphi(\bm{r})\big)\,\kappa(x),
\]
where $\bm{a},\bm{\theta}$ are $x$–dependent translations/rotations and $\varphi:\Omega\to\mathbb{R}$ is a \emph{warping displacement} along $\mathbf{i}$ with scalar amplitude $\kappa(x)$. 
The Saint–Venant class implies that the sectional shear strain vector equals the in–plane gradient of the axial warping:
\[
\bm{\gamma}(\bm{r})=2\,\bm{\varepsilon}(\bm{u})\,\mathbf{i}\ -\ \big(\mathbf{i}\cdot 2\,\bm{\varepsilon}(\bm{u})\,\mathbf{i}\big)\,\mathbf{i}
\;\;=\;\nabla_{\bm{r}}\big(\kappa(x)\,\varphi(\bm{r})\big).
\]
Since $\bm{\tau}=2G\,\bm{\varepsilon}\,\mathbf{i}-\sigma\,\mathbf{i}$ projects in the section to $\bm{\tau}=2G\,\bm{\gamma}$, we obtain the kinematic identification
\begin{equation}
\bm{\tau}(\bm{r})=2G\,\nabla_{\bm{r}}\big(\kappa\,\varphi(\bm{r})\big),
\qquad 
\bm{\tau}\cdot\bm{\nu}=0\ \Rightarrow\ \partial_{\bm{\nu}}\big(\kappa\,\varphi\big)=0\ \text{on }\partial\Omega.
\label{eq:tau-from-warp}
\end{equation}
Eliminating $\kappa$ into the global force constraint in \eqref{eq:resultants} fixes the amplitude (up to rigid motions of $\bm{u}$), while $\varphi$ is determined (up to an additive constant) by the plane Neumann problem that is \emph{equivalent} to \eqref{eq:SV-variational}:
\begin{equation}
\text{find }\varphi\in H^1(\Omega)/\mathbb{R}\ \text{s.t.}\quad
\int_\Omega \nabla_{\bm{r}}\varphi\cdot\nabla_{\bm{r}}\eta\,dA
=\frac{1}{2G\,\kappa}\,\bm{F}\cdot\int_{\partial\Omega}\eta\,\bm{\nu}\,ds,
\quad\forall\eta\in H^1(\Omega),
\label{eq:phi-weak}
\end{equation}
together with the torsion–free side constraint
\begin{equation}
\int_\Omega \bm{r}\times\nabla_{\bm{r}}\big(\kappa\,\varphi\big)\,dA=\mathbf{0}.
\label{eq:phi-tf}
\end{equation}
Equations \eqref{eq:tau-from-warp}–\eqref{eq:phi-weak} \emph{define} the warping formulation without recourse to a stress potential; the resulting $\bm{\tau}$ then generates $\sigma$ by \eqref{eq:sigma-from-tau}.

\subsection*{Solution summary}

Given $\bm{F}=F\,\bm{f}\in\mathbf{i}^\perp$:

\begin{enumerate}
\item Solve the plane constrained problem \eqref{eq:SV-variational} (or, equivalently, \eqref{eq:phi-weak}–\eqref{eq:phi-tf}) on $\Omega$ to obtain the unique shear field $\bm{\tau}(\bm{r})$ with $\bm{\tau}\cdot\bm{\nu}=0$ on $\partial\Omega$, $\int_\Omega\bm{\tau}=\bm{F}$, and $\int_\Omega\bm{r}\times\bm{\tau}=\mathbf{0}$.

\item Reconstruct axial stress by $\sigma(x,\bm{r})=(L-x)\,\operatorname{div}_{\bm{r}}\bm{\tau}(\bm{r})$.

\item The full Cauchy stress is
\[
\bm{T}(x,\bm{r})=\sigma(x,\bm{r})\,\mathbf{i}\otimes\mathbf{i}
+\bm{\tau}(\bm{r})\otimes\mathbf{i}
+\mathbf{i}\otimes\bm{\tau}(\bm{r}),
\]
which satisfies equilibrium, traction–free sides, the end resultants, and the torsion–free condition.

\item The bending moment vector satisfies $\bm{M}(x)=(L-x)\,\mathbf{i}\times\bm{F}$ and $\bm{V}=\bm{F}$, cf. \eqref{eq:moment-balance}.
\end{enumerate}

\paragraph{Remarks.}
(i) The representation \eqref{eq:frame-stress-tensor} is the most general symmetric tensor compatible with traction–free lateral boundary under the Saint–Venant class; \eqref{eq:side-BC} is necessary and sufficient for side tractions to vanish.
(ii) The formulation does not assume any special cross–section; all section dependence is carried by the plane problem for $\bm{\tau}$ (or, kinematically, by $\varphi$).
(iii) Uniqueness holds up to the usual self‐equilibrated rigid modes: adding a self‐equilibrated $(\sigma,\bm{\tau})$ with vanishing resultants leaves the solution class invariant; \eqref{eq:SV-variational} singles out the Saint–Venant solution by minimum complementary energy.
